\documentclass[12pt]{article}
\usepackage[utf8]{inputenc}
\usepackage[T1]{fontenc}
\usepackage{hyperref}
\title{Ukwa Language Resources for K-12 Curriculum Development}
\author{Compiled by Claude Sonnet 5 (claude-sonnet-5) for Andr\'as Kornai}
\date{}
\begin{document}
\maketitle

\section{Introduction}
Ukwa (ISO 639-3 ukq; Glottocode ukwa1238) is an Ibibio-Efik language of Cross River State, Nigeria. Sources found in search disagree sharply on the number of speakers: one Wikipedia mirror gives about 100 (2004), while the current English Wikipedia infobox gives 350,000 (2023). Joshua Project reports an ethnic population of about 700, mostly in Odukpani Local Government Area. Do not confuse it with Ukwa West, a Local Government Area in Abia State whose people speak an Igbo dialect (Asa). Very little educational material was found.

\section{Primary Educational Resources}
No primers, textbooks or school programmes for Ukwa were found in search results. Joshua Project states that no gospel materials exist in Ukwa itself.

\section{Linguistic and Reference Materials}
\begin{itemize}
\item Wikipedia, ``Ukwa language'': \url{https://en.wikipedia.org/wiki/Ukwa_language}
\item Wikipedia mirror, ISO 639:ukq (older speaker count): \url{https://troywiki.duckdns.org/content/wikipedia_en_all_maxi_2024-01/A/ISO_639%3Aukq}
\item Wikipedia, ``Ukwa (disambiguation)'': \url{https://en.wikipedia.org/wiki/Ukwa_(disambiguation)}
\end{itemize}

\section{Religious and Community Resources}
\begin{itemize}
\item Joshua Project, Ukwa in Nigeria (people group profile): \url{https://joshuaproject.net/people_groups/18868}
\item Joshua Project language page for ukq: \url{https://joshuaproject.net/languages/ukq}
\item Scripture Earth page for ukq, which links to a Ukwa Bible on YouVersion (not verified here): \url{https://www.scriptureearth.org/00eng.php?iso=ukq}
\end{itemize}

\section{Digital and Audio Resources}
No audio, video or software resources were found. Joshua Project shows no Audio Bible or Jesus Film.

\section{Curriculum Development Context}
The language is not known to be taught in schools, and the evidence found is thin and partly contradictory. A curriculum would need new fieldwork to settle speaker numbers, orthography and community priorities. The related and much better documented Efik and Ibibio may offer materials for adaptation, but none were searched here.

\end{document}